% Part: proof-theory
% Chapter: normalization
% Section: permutations

\documentclass[../../../include/open-logic-section]{subfiles}

\begin{document}

\olfileid{pt}{nor}{per}

\olsection{تبدیل‌های جایگشتی}

برش‌ها در \Log{N1i} قطعه‌هایی هستند که به مقدمهٔ عمدهٔ
!!a{elimination} پایان می‌یابند. یکی از حالت‌های اثبات نرمال‌سازی، کاهش
طول برش‌های بیشینه و همراه با آن کاهش طول برشِ !!{proof} است. برشی با
طول $>1$ ممکن است به صورت‌های گوناگونی پایان یابد، بسته به اینکه آخرین
رخداد !!{formula}~$!A$ در برش نتیجهٔ \Elim{\lor} یا
\Elim{\lexists} باشد و بسته به اینکه مقدمهٔ عمدهٔ کدام استنتاج
!!{elimination} باشد.

برای نمونه، اگر استنتاج‌های درگیر \Elim{\lor} و \Elim{\land} باشند،
فرمول برش $!A\ident !B\land !C$ است و زیر-!!{proof}~$\delta_1$ای داریم
که چنین پایان می‌یابد:
\[
\AxiomC{}
\RightLabel{$\delta_2$}
\DeduceC{$!D \lor !E$}
\AxiomC{$\Discharge{!D}{x}$}
\RightLabel{$\delta_3$}
\DeduceC{$!B \land !C$}
\AxiomC{$\Discharge{!E}{y}$}
\RightLabel{$\delta_4$}
\DeduceC{$!B \land !C$}
\DischargeRule{\Elim{\lor}}{x\,y}
\TrinaryInfC{$!B \land !C$}
\RightLabel{\Elim{\land}[1]}
\UnaryInfC{$!B$}
\DisplayProof
\]
آخرین رخداد !!{formula}~$!A_n$ در برش، نتیجهٔ \Elim{\lor}، یعنی
$!B\land !C$ پایینی، است. رخداد یکی‌مانده‌به‌آخر !!{formula}
~$!A_{n-1}$ یکی از مقدمات فرعی \Elim{\lor} است.

برای کاهش طول این برش می‌توانیم ترتیب استنتاج‌های \Elim{\lor} و
\Elim{\land} را عوض، یا «جایگشت»، کنیم؛ یعنی زیر‌برهان~$\delta_1$ در
~‌$\delta$ را با برهان زیر جایگزین کنیم:
\[
\AxiomC{}
\RightLabel{$\delta_2$}
\DeduceC{$!D \lor !E$}
\AxiomC{$\Discharge{!D}{x}$}
\RightLabel{$\delta_3$}
\DeduceC{$!B \land !C$}
\RightLabel{\Elim{\land}[1]}
\UnaryInfC{$!B$}
\AxiomC{$\Discharge{!E}{y}$}
\RightLabel{$\delta_4$}
\DeduceC{$!B \land !C$}
\RightLabel{\Elim{\land}[1]}
\UnaryInfC{$!B$}
\DischargeRule{\Elim{\lor}}{x\,y}
\TrinaryInfC{$!B$}
\DisplayProof
\]
برش‌هایی که پیش‌تر در مقدمهٔ \Elim{\land} زیر \Elim{\lor} پایان
می‌یافتند، اکنون در مقدمهٔ یکی از استنتاج‌های \Elim{\land} بالای
مقدمات فرعی \Elim{\lor} پایان می‌یابند؛ پس طول هر یک یک واحد کاهش
یافته است.

افزون بر این، هر برشی که سراسر در $\delta_2$، $\delta_3$ یا
~‌$\delta_4$، یا سراسر بیرون از~$\delta_1$ در~$\delta$ قرار دارد
بی‌تغییر می‌ماند: همان !!{formula}ها را در بر می‌گیرد و از همان
استنتاج‌ها می‌گذرد، پس به‌ویژه رتبه و طول آن با برش متناظر در~$\delta$
یکسان است. بنابراین رتبهٔ برشِ !!{proof} تازه افزایش نیافته و طول برشِ
!!{proof} تازه از برهان پیشین کمتر است.

\begin{prob}
پس از یک تبدیل جایگشتی، طول دست‌کم یک برش یک واحد کم می‌شود. اگر
\Elim{\land} را از میان \Elim{\lor} جایگشت کنیم، در واقع طول دو قطعهٔ
برشی کاهش می‌یابد؛ پس در این حالت طول برشِ !!{proof} دست‌کم دو واحد
کم می‌شود. این نتیجه برای هر !!{proof} برقرار است. آیا طول برشِ
!!a{proof} ممکن است با یک چنین جایگشتی بیش از
دو واحد کاهش یابد؟ آیا \emph{رتبهٔ} برش می‌تواند کاهش یابد؟
\end{prob}

فرایند انتقال !!a{elimination} به بالای قاعدهٔ \Elim{\lor} یا
\Elim{\lexists} را \emph{تبدیل جایگشتی} می‌نامند. الگوهای برخی
جفت‌قاعده‌ها در \olref{tab:permutation-conversions} آمده‌اند و بقیه به
عنوان تمرین واگذار می‌شوند.

% \begin{table}
% \begin{multline*}
% \shoveleft{\AxiomC{$!D \lor !E$}
% \AxiomC{$\Discharge{!D}{x}$}
% \shortDeduce
% \DeduceC{$!B \land !C$}
% \AxiomC{$\Discharge{!E}{y}$}
% \shortDeduce
% \DeduceC{$!B \land !C$}
% \DischargeRule{\Elim{\lor}}{x\,y}
% \TrinaryInfC{$!B \land !C$}
% \RightLabel{\Elim{\land}[1]}
% \UnaryInfC{$!B$}
% \DisplayProof}\\
% \shoveright{\longrightarrow\ 
% \AxiomC{$!D \lor !E$}
% \AxiomC{$\Discharge{!D}{x}$}
% \shortDeduce
% \DeduceC{$!B \land !C$}
% \RightLabel{\Elim{\land}[1]}
% \UnaryInfC{$!B$}
% \AxiomC{$\Discharge{!E}{y}$}
% \shortDeduce
% \DeduceC{$!B \land !C$}
% \RightLabel{\Elim{\land}[1]}
% \UnaryInfC{$!B$}
% \DischargeRule{\Elim{\lor}}{x\,y}
% \TrinaryInfC{$!B$}
% \DisplayProof}
% \\
% \shoveleft{\AxiomC{$!D \lor !E$}
% \AxiomC{$\Discharge{!D}{x}$}
% \shortDeduce
% \DeduceC{$!B \lif !C$}
% \AxiomC{$\Discharge{!E}{y}$}
% \shortDeduce
% \DeduceC{$!B \lif !C$}
% \DischargeRule{\Elim{\lor}}{x\,y}
% \TrinaryInfC{$!B \lif !C$}
% \AxiomC{$!B$}
% \RightLabel{\Elim{\lif}}
% \BinaryInfC{$!C$}
% \DisplayProof}\\
% \shoveright{\longrightarrow\ 
% \AxiomC{$!D \lor !E$}
% \AxiomC{$\Discharge{!D}{x}$}
% \shortDeduce
% \DeduceC{$!B \lif !C$}
% \AxiomC{$!B$}
% \RightLabel{\Elim{\lif}}
% \BinaryInfC{$!C$}
% \AxiomC{$\Discharge{!E}{y}$}
% \shortDeduce
% \DeduceC{$!B \lif !C$}
% \AxiomC{$!B$}
% \RightLabel{\Elim{\lif}}
% \BinaryInfC{$!C$}
% \DischargeRule{\Elim{\lor}}{x\,y}
% \TrinaryInfC{$!C$}
% \DisplayProof}
% \\
% \shoveleft{\AxiomC{$!D \lor !E$}
% \AxiomC{$\Discharge{!D}{x}$}
% \shortDeduce
% \DeduceC{$!B \lif !C$}
% \AxiomC{$\Discharge{!E}{y}$}
% \shortDeduce
% \DeduceC{$!B \lif !C$}
% \DischargeRule{\Elim{\lor}}{x\,y}
% \TrinaryInfC{$!B \lif !C$}
% \AxiomC{$!B$}
% \RightLabel{\Elim{\lif}}
% \BinaryInfC{$!C$}
% \DisplayProof}\\
% \shoveright{\longrightarrow\ 
% \AxiomC{$!D \lor !E$}
% \AxiomC{$\Discharge{!D}{x}$}
% \shortDeduce
% \DeduceC{$!B \lif !C$}
% \AxiomC{$!B$}
% \RightLabel{\Elim{\lif}}
% \BinaryInfC{$!C$}
% \AxiomC{$\Discharge{!E}{y}$}
% \shortDeduce
% \DeduceC{$!B \lif !C$}
% \AxiomC{$!B$}
% \RightLabel{\Elim{\lif}}
% \BinaryInfC{$!C$}
% \DischargeRule{\Elim{\lor}}{x\,y}
% \TrinaryInfC{$!C$}
% \DisplayProof}
% \\
% \shoveleft{\AxiomC{$!D \lor !E$}
% \AxiomC{$\Discharge{!D}{x}$}
% \shortDeduce
% \DeduceC{$\lforall[x][!B(x)]$}
% \AxiomC{$\Discharge{!E}{y}$}
% \shortDeduce
% \DeduceC{$\lforall[x][!B(x)]$}
% \DischargeRule{\Elim{\lor}}{x\,y}
% \TrinaryInfC{$\lforall[x][!B(x)]$}
% \RightLabel{\Elim{\lforall}}
% \UnaryInfC{$!B(t)$}
% \DisplayProof}\\
% \shoveright{\longrightarrow\ 
% \AxiomC{$!D \lor !E$}
% \AxiomC{$\Discharge{!D}{x}$}
% \shortDeduce
% \DeduceC{$\lforall[x][!B(x)]$}
% \RightLabel{\Elim{\lforall}}
% \UnaryInfC{$!B(t)$}
% \AxiomC{$\Discharge{!E}{y}$}
% \shortDeduce
% \DeduceC{$\lforall[x][!B(x)]$}
% \RightLabel{\Elim{\lforall}}
% \UnaryInfC{$!B(t)$}
% \DischargeRule{\Elim{\lor}}{x\,y}
% \TrinaryInfC{$!B(t)$}
% \DisplayProof}
% \end{multline*}
% \caption{Some permutation conversions for \Log{N1i}.}
% \ollabel{tab:permutation-conversions}
% \end{table}

{\small
\begin{longtable}{@{}l@{}}
\AxiomC{$!D \lor !E$}
\AxiomC{$\Discharge{!D}{x}$}
\shortDeduce
\DeduceC{$!B \land !C$}
\AxiomC{$\Discharge{!E}{y}$}
\shortDeduce
\DeduceC{$!B \land !C$}
\DischargeRule{\Elim{\lor}}{x\,y}
\TrinaryInfC{$!B \land !C$}
\RightLabel{\Elim{\land}[1]}
\UnaryInfC{$!B$}
\DisplayProof
$\longrightarrow$\\[6ex]
\quad\AxiomC{$!D \lor !E$}
\AxiomC{$\Discharge{!D}{x}$}
\shortDeduce
\DeduceC{$!B \land !C$}
\RightLabel{\Elim{\land}[1]}
\UnaryInfC{$!B$}
\AxiomC{$\Discharge{!E}{y}$}
\shortDeduce
\DeduceC{$!B \land !C$}
\RightLabel{\Elim{\land}[1]}
\UnaryInfC{$!B$}
\DischargeRule{\Elim{\lor}}{x\,y}
\TrinaryInfC{$!B$}
\DisplayProof
\\[10ex]
\AxiomC{$!D \lor !E$}
\AxiomC{$\Discharge{!D}{x}$}
\shortDeduce
\DeduceC{$!B \lif !C$}
\AxiomC{$\Discharge{!E}{y}$}
\shortDeduce
\DeduceC{$!B \lif !C$}
\DischargeRule{\Elim{\lor}}{x\,y}
\TrinaryInfC{$!B \lif !C$}
\AxiomC{$!B$}
\RightLabel{\Elim{\lif}}
\BinaryInfC{$!C$}
\DisplayProof
$\longrightarrow$\\[6ex]
\AxiomC{$!D \lor !E$}
\AxiomC{$\Discharge{!D}{x}$}
\shortDeduce
\DeduceC{$!B \lif !C$}
\AxiomC{$!B$}
\RightLabel{\Elim{\lif}}
\BinaryInfC{$!C$}
\AxiomC{$\Discharge{!E}{y}$}
\shortDeduce
\DeduceC{$!B \lif !C$}
\AxiomC{$!B$}
\RightLabel{\Elim{\lif}}
\BinaryInfC{$!C$}
\DischargeRule{\Elim{\lor}}{x\,y}
\TrinaryInfC{$!C$}
\DisplayProof
\\[10ex]
\AxiomC{$!D \lor !E$}
\AxiomC{$\Discharge{!D}{x}$}
\shortDeduce
\DeduceC{$\lforall[x][!B(x)]$}
\AxiomC{$\Discharge{!E}{y}$}
\shortDeduce
\DeduceC{$\lforall[x][!B(x)]$}
\DischargeRule{\Elim{\lor}}{x\,y}
\TrinaryInfC{$\lforall[x][!B(x)]$}
\RightLabel{\Elim{\lforall}}
\UnaryInfC{$!B(t)$}
\DisplayProof
\quad$\longrightarrow$\\[8ex]
\quad\AxiomC{$!D \lor !E$}
\AxiomC{$\Discharge{!D}{x}$}
\shortDeduce
\DeduceC{$\lforall[x][!B(x)]$}
\RightLabel{\Elim{\lforall}}
\UnaryInfC{$!B(t)$}
\AxiomC{$\Discharge{!E}{y}$}
\shortDeduce
\DeduceC{$\lforall[x][!B(x)]$}
\RightLabel{\Elim{\lforall}}
\UnaryInfC{$!B(t)$}
\DischargeRule{\Elim{\lor}}{x\,y}
\TrinaryInfC{$!B(t)$}
\DisplayProof
\\
\caption{چند تبدیل جایگشتی برای \Log{N1i}.}
\ollabel{tab:permutation-conversions}
\end{longtable}
}

\begin{prob}
تبدیل‌های جایگشتی را برای \Elim{\lor} که \Elim{\lor} یا
\Elim{\lnot} در پی آن می‌آید بنویسید.
\end{prob}

\begin{prob}
تبدیل‌های جایگشتی را برای \Elim{\lexists} که \Elim{\lexists} در پی
آن می‌آید بنویسید. توضیح دهید چرا در حالت دوم هیچ شرط متغیر ویژه‌ای
نقض نمی‌شود.
\end{prob}

باید مطمئن شویم که اعمال تبدیل جایگشتی ممکن نیست شرط متغیر ویژه‌ای را
نقض کند. برهان زیر را در نظر بگیرید:
\[
\AxiomC{}
\RightLabel{$\delta_2$}
\DeduceC{$!D \lor !E$}
\AxiomC{$\Discharge{!D}{x}$}
\RightLabel{$\delta_3$}
\DeduceC{$\lexists[x][!B(x)]$}
\AxiomC{$\Discharge{!E}{y}$}
\RightLabel{$\delta_4$}
\DeduceC{$\lexists[x][!B(x)]$}
\DischargeRule{\Elim{\lor}}{x\,y}
\TrinaryInfC{$\lexists[x][!B(x)]$}
\AxiomC{$\Discharge{!B(c)}{z}$}
\RightLabel{$\delta_5$}
\DeduceC{$!C$}
\DischargeRule{\Elim{\lexists}}{z}
\BinaryInfC{$!C$}
\DisplayProof
\]
اگر آن را با برهان زیر جایگزین کنیم، دو نسخهٔ~$\delta_5$ را نخست با
متغیرهای ویژه و برچسب‌های ترخیص دوبه‌دو تازه تغییر نام می‌دهیم:
\[
\AxiomC{}
\RightLabel{$\delta_2$}
\DeduceC{$!D \lor !E$}
\AxiomC{$\Discharge{!D}{x}$}
\RightLabel{$\delta_3$}
\DeduceC{$\lexists[x][!B(x)]$}
\AxiomC{$\Discharge{!B(c_1)}{z_1}$}
\RightLabel{$\delta_5^{(1)}$}
\DeduceC{$!C$}
\DischargeRule{\Elim{\lexists}}{z_1}
\BinaryInfC{$!C$}
\AxiomC{$\Discharge{!E}{y}$}
\RightLabel{$\delta_4$}
\DeduceC{$\lexists[x][!B(x)]$}
\AxiomC{$\Discharge{!B(c_2)}{z_2}$}
\RightLabel{$\delta_5^{(2)}$}
\DeduceC{$!C$}
\DischargeRule{\Elim{\lexists}}{z_2}
\BinaryInfC{$!C$}
\DischargeRule{\Elim{\lor}}{x\,y}
\TrinaryInfC{$!C$}
\DisplayProof
\]
ممکن است نگران باشیم که متغیر ویژهٔ~$c_i$، برای نمونه، در~$!D$ ظاهر
شود. در برهان اصلی، فرض~$!D$ بالای \Elim{\lexists} مرخص می‌شود و در
فرضی که بالای مقدمهٔ عمدهٔ استنتاج \Elim{\lexists} باز است ظاهر
نمی‌شود؛ پس شرط متغیر ویژه آنجا برقرار است. در !!{proof} تازه، $!D$ تا
پس از استنتاج‌های \Elim{\lexists} مرخص نمی‌شود و بدون تغییر نام شرط
متغیر ویژه ممکن بود نقض شود. اما به یاد آورید که !!{proof}های ما منظم
و نسخه‌های~$\delta_5$ با متغیرهای ویژهٔ دوبه‌دو تازه‌اند: هر متغیر ویژه
فقط متعلق به یک استنتاج است و تنها بالای مقدمهٔ فرعی \Elim{\lexists}
در !!{proof} متناظر ظاهر می‌شود. به‌ویژه، در همین !!{proof} هیچ‌یک از
$c_1,c_2$ در
$\delta_3$ یا~$\delta_4$ ظاهر نمی‌شود.

این مثال نکتهٔ دیگری را نیز نشان می‌دهد: !!{proof} تازه دو نسخه از
~‌$\delta_5$ دارد. اگر $\delta_5$ دارای برشی بیشینه باشد، این تکثیر
ممکن است طول برش را ثابت نگه دارد یا حتی افزایش دهد. پس باید تبدیل
جایگشتی را فقط هنگامی اعمال کنیم که بدانیم هیچ برش بیشینه‌ای در
~‌$\delta_5$ نیست. این امر را با انتخاب همیشگی یک برش بیشینهٔ
\emph{بالاترین} تضمین می‌کنیم؛ یعنی برشی بیشینه که زیر-!!{proof} آن---
اینجا~$\delta_1$---هیچ برش بیشینه‌ای ندارد که بالای آخرین استنتاج در
زیر-!!{proof} پایان یابد. بدین‌ترتیب امکان وجود برش‌های بیشینهٔ دیگر در
~$\delta_5$، و در واقع در $\delta_2$، $\delta_3$ یا~$\delta_4$، منتفی
می‌شود.

\end{document}
